% !TEX root = ../main.tex
\chapter{命題・証明・仕様}
\label{chap:ch02_props_specs}

\begin{chapterabstract}
本章では、Lean 4 における命題と証明を、ソフトウェア仕様を検査可能な形にするための道具として導入する。
多くのプログラミング言語では、条件式は \texttt{true} または \texttt{false} を返す値として扱われる。
しかし仕様を扱うときには、「この関数はすべての入力でこの性質を満たす」「この変換は ID を保存する」のような一般的な主張を書きたくなる。
Lean では、そのような主張を \texttt{Prop} として表し、\texttt{example} や \texttt{theorem} によって証明する。
本章の目標は、証明の専門技術を網羅することではない。
関数に対して「何を保証したいのか」を命題として切り出し、\texttt{rfl}、\texttt{simp}、\texttt{rw} で小さく確認する感覚を得ることである。
\end{chapterabstract}

\begin{learninggoals}
\item \texttt{Bool} と \texttt{Prop} の違いを、実行時の判定と仕様上の主張の違いとして説明できる。
\item 関数が満たすべき性質を、Lean の命題として書ける。
\item \texttt{example}、\texttt{theorem}、\texttt{by} を使って小さな証明を書ける。
\item \texttt{rfl}、\texttt{simp}、\texttt{rw} の役割を、仕様確認の道具として読み分けられる。
\item 税込金額計算と ID 保存のような小さな実務仕様を、Lean の定理として表せる。
\end{learninggoals}

\section{この章を読む理由}

第1章では、Lean で型、値、関数を書く方法を見た。
たとえば、価格に税額を足す関数、入力を正規化する関数、データを別の形へ変換する関数は、\texttt{def} で定義できる。
しかし、関数が書けるだけでは仕様を扱ったことにはならない。
実務で本当に知りたいのは、しばしば次のような性質である。

\begin{itemize}
\item 税額が 0 のとき、税込金額は元の価格と同じになる。
\item DTO 変換をしても、ユーザー ID は変わらない。
\item 会員向けの送料計算では、会員の送料は必ず 0 になる。
\item 正規化関数を通してから計算しても、結果は変わらない。
\end{itemize}

これらは、1回の \texttt{\#eval} で確認できる例ではない。
「この入力では合っている」ではなく、「すべての入力でこの形の性質が成り立つ」と言いたいからである。

\begin{intuitionbox}
\texttt{Bool} は、プログラムが計算して返す真偽値である。
\texttt{Prop} は、Lean に対して「これを証明したい」と提示する主張である。
仕様を書くときには、実行時の判定だけでなく、証明対象としての主張を区別する必要がある。
\end{intuitionbox}

\FigurePlaceholder{fig-ch02-overview.png}{0.90}{第2章の概念地図：関数、命題、証明、仕様の関係}{fig:ch02-overview}

図\ref{fig:ch02-overview}は、本章の流れを表す。
まず関数がある。
次に、その関数に期待する性質を命題として書く。
最後に、その命題が正しいことを証明として与える。
この流れが、後続章で扱うリファクタリング証明、マイグレーション証明、API adapter の一貫性確認の土台になる。

\section{Bool と Prop の違い}

\subsection{Bool は計算される値である}

\texttt{Bool} は、\texttt{true} または \texttt{false} という値を持つ型である。
多くのプログラミング言語の boolean とほぼ同じように読める。
たとえば、会員なら割引対象にする、一定金額以上なら送料無料にする、といった分岐で使う。

\begin{listing}[htbp]
\begin{minted}{lean4}
def hasDiscount (member : Bool) : Bool :=
  member

#eval hasDiscount true   -- true
#eval hasDiscount false  -- false
\end{minted}
\caption{Bool は計算される真偽値である}
\label{lst:ch02_bool}
\end{listing}

コード\ref{lst:ch02_bool}では、\texttt{hasDiscount true} を評価すると \texttt{true} が得られる。
これは、実行時に計算される値である。
テストコードで特定の入力を与えて結果を確認する感覚に近い。

\begin{softwarebox}
\texttt{Bool} は、アプリケーション内の分岐や判定結果に向いている。
たとえば「このリクエストは有効か」「このユーザーは会員か」「この金額は送料無料条件を満たすか」を実行時に判定するときに使う。
\end{softwarebox}

\subsection{Prop は証明したい主張である}

一方、\texttt{Prop} は「証明対象となる命題」の型である。
値として計算することよりも、その主張が正しいことを Lean に納得させることが目的になる。

\begin{listing}[htbp]
\begin{minted}{lean4}
def DiscountSpec (member : Bool) : Prop :=
  hasDiscount member = true

theorem member_hasDiscount : DiscountSpec true := by
  rfl
\end{minted}
\caption{Prop は証明したい主張である}
\label{lst:ch02_prop}
\end{listing}

コード\ref{lst:ch02_prop}の \texttt{DiscountSpec true} は、\texttt{hasDiscount true = true} という主張である。
\texttt{theorem member\_hasDiscount} は、その主張が正しいことを Lean に示している。
ここで使っている \texttt{rfl} は、左右が定義を展開すれば同じになることを示す証明である。

\begin{definitionbox}
本書では、\textbf{命題}を「正しいかどうかを証明したい主張」として扱う。
Lean では命題の型を \texttt{Prop} と書く。
たとえば \texttt{price = addTax base tax} や \texttt{(toDto u).id = u.id} は命題である。
\end{definitionbox}

\subsection{仕様では両方を使う}

\texttt{Bool} と \texttt{Prop} は競合するものではない。
実務仕様では、しばしば両方を使う。
\texttt{Bool} は実装で判定する値として使い、その判定関数が期待どおりに振る舞うことを \texttt{Prop} として表す。

\begin{table}[htbp]
\centering
\caption{\texttt{Bool} と \texttt{Prop} の使い分け}
\label{tab:ch02-bool-prop}
\begin{tabular}{p{0.22\linewidth}p{0.34\linewidth}p{0.34\linewidth}}
\toprule
観点 & \texttt{Bool} & \texttt{Prop} \\
\midrule
役割 & 実行時に計算される真偽値 & 証明したい主張 \\
主な使い道 & 条件分岐、バリデーション、フラグ & 仕様、不変条件、事後条件 \\
確認方法 & \texttt{\#eval} やテストで値を見る & \texttt{theorem} や \texttt{example} で証明する \\
例 & \texttt{hasDiscount member} & \texttt{hasDiscount member = true} \\
\bottomrule
\end{tabular}
\end{table}

表\ref{tab:ch02-bool-prop}の違いを、最初から完全に使いこなす必要はない。
重要なのは、\texttt{Bool} は「判定結果」、\texttt{Prop} は「保証したい主張」と読むことである。

\FigurePlaceholder{fig-ch02-bool-vs-prop.png}{0.88}{\texttt{Bool} と \texttt{Prop} の違い：実行時の判定と仕様上の主張}{fig:ch02-bool-vs-prop}

\section{命題としての仕様}

\subsection{自然言語仕様を小さく切り出す}

自然言語の仕様書には、さまざまな粒度の文が混ざっている。
たとえば、次のような仕様を考える。

\begin{quote}
税込金額は、税抜価格に税額を加えた金額である。
税額が 0 の場合、税込金額は税抜価格と一致する。
\end{quote}

この文章をいきなりすべて Lean に移す必要はない。
まずは、関数と命題に分ける。

\begin{itemize}
\item 関数：税抜価格と税額から税込金額を計算する。
\item 命題：税額が 0 のとき、結果は元の価格と等しい。
\end{itemize}

\begin{listing}[htbp]
\begin{minted}{lean4}
def addTax (price tax : Nat) : Nat :=
  price + tax

def TaxSpec (price tax total : Nat) : Prop :=
  total = addTax price tax

example : TaxSpec 1000 100 1100 := by
  rfl
\end{minted}
\caption{仕様を関数と命題に分ける}
\label{lst:ch02_spec_as_prop}
\end{listing}

コード\ref{lst:ch02_spec_as_prop}では、\texttt{addTax} が計算であり、\texttt{TaxSpec} が仕様である。
\texttt{TaxSpec 1000 100 1100} は「1000 に 100 を足した税込金額は 1100 である」という命題になる。

\begin{leanbox}
Lean では、仕様をコメントとしてではなく、型のついた命題として書ける。
その命題に証明を与えると、Lean が主張と根拠の対応を検査する。
この点が、通常のコメントや設計メモとの大きな違いである。
\end{leanbox}

\subsection{仕様は関数の外側に置く}

実装関数と仕様命題を分けると、後で変更に強くなる。
関数の中に「こういうつもり」とコメントを書くだけでは、Lean はそれを検査できない。
しかし命題として外に出せば、関数を変更したときに証明が壊れる可能性がある。
壊れた証明は、仕様上の意味が変わったかもしれないというシグナルになる。

\begin{softwarebox}
実務では、すべての仕様を Lean に移す必要はない。
まずは「壊れると困る核」を選ぶ。
たとえば、ID が保存されること、件数が保存されること、手数料が 0 のとき結果が変わらないこと、round-trip が成り立つことなどである。
\end{softwarebox}

\section{example と theorem}

\subsection{example は名前のない小さな確認である}

\texttt{example} は、名前を付けずに命題と証明を書くための構文である。
本文中の小さな確認や、Lean の構文を試すときに向いている。

\begin{listing}[htbp]
\begin{minted}{lean4}
def SamePrice (before after : Nat) : Prop :=
  before = after

example : SamePrice 1200 1200 := by
  rfl
\end{minted}
\caption{example による小さな証明}
\label{lst:ch02_example}
\end{listing}

コード\ref{lst:ch02_example}は、「1200 と 1200 は同じ価格である」という命題を証明している。
この証明に名前はない。
そのため、後から別の証明で参照することはできない。

\subsection{theorem は名前付きの仕様である}

\texttt{theorem} は、名前付きの命題と証明を書くための構文である。
後から参照したい仕様、設計文書に対応させたい性質、CI で壊れたときに意味を追跡したい性質には、\texttt{theorem} を使う。

\begin{listing}[htbp]
\begin{minted}{lean4}
theorem addTax_zero (price : Nat) :
    addTax price 0 = price := by
  simp [addTax]
\end{minted}
\caption{theorem による名前付き仕様}
\label{lst:ch02_theorem}
\end{listing}

コード\ref{lst:ch02_theorem}は、「税額が 0 のとき、税込金額は元の価格と一致する」という名前付きの仕様である。
\texttt{addTax\_zero} という名前が付いているため、後の章でこの事実を使って別の証明を組み立てることができる。

\begin{definitionbox}
本書では、\texttt{example} を「その場限りの確認」、\texttt{theorem} を「名前を付けて再利用したい仕様」として使い分ける。
実務の設計文書と対応させるなら、重要な性質には \texttt{theorem} 名を付けるとよい。
\end{definitionbox}

\section{by で証明を書く}

\subsection{by は証明モードへの入口である}

Lean では、命題の後に \texttt{:= by} と書くと、証明を書くモードに入る。
ここから先では、\texttt{rfl}、\texttt{simp}、\texttt{rw}、\texttt{intro} などの証明コマンドを使って、ゴールを解消していく。

\begin{listing}[htbp]
\begin{minted}{lean4}
theorem same_price_refl (price : Nat) :
    price = price := by
  rfl
\end{minted}
\caption{by ブロックの基本形}
\label{lst:ch02_by_proof}
\end{listing}

コード\ref{lst:ch02_by_proof}では、\texttt{price = price} という命題を証明している。
これはどの \texttt{price} についても明らかに成り立つ。
Lean では、その明らかさを \texttt{rfl} という証明で伝える。

\subsection{証明はコメントではない}

証明は、コードに添える説明文ではない。
Lean が検査する構造化された根拠である。
たとえば、コメントなら間違っていてもコンパイルは通る。
しかし theorem の証明が間違っていれば、Lean は受け付けない。

\begin{warningbox}
\texttt{theorem} の名前がもっともらしくても、証明がなければ保証にはならない。
本書では原則として \texttt{sorry} を使わない。
\texttt{sorry} は「ここは未証明だが、いったん通す」という印であり、仕様保証としては未完了である。
\end{warningbox}

\section{rfl、simp、rw の最初の使い方}

本章では、証明の道具を三つだけに絞って導入する。
\texttt{rfl}、\texttt{simp}、\texttt{rw} である。
これらだけでも、小さな仕様確認には十分役に立つ。

\subsection{rfl は定義を展開すれば同じことを示す}

\texttt{rfl} は reflexivity、つまり反射律に由来する。
雑に言えば、左右が定義を展開すると同じものになるときに使える。

\begin{listing}[htbp]
\begin{minted}{lean4}
theorem addTax_def (price tax : Nat) :
    addTax price tax = price + tax := by
  rfl
\end{minted}
\caption{rfl で定義どおりの等式を証明する}
\label{lst:ch02_rfl}
\end{listing}

コード\ref{lst:ch02_rfl}では、\texttt{addTax price tax} の定義が \texttt{price + tax} なので、\texttt{rfl} で証明できる。
これは「定義どおりです」と Lean に伝えていると読める。

\subsection{simp は明らかな整理を任せる}

\texttt{simp} は、Lean が知っている単純化規則を使って式を整理する道具である。
たとえば、\texttt{x + 0 = x} のような自然数の基本的な整理や、\texttt{if true then ... else ...} のような分岐の整理に使える。

\begin{listing}[htbp]
\begin{minted}{lean4}
def chargeShipping (isMember : Bool) (shipping : Nat) : Nat :=
  if isMember then 0 else shipping

theorem member_shipping_zero (shipping : Nat) :
    chargeShipping true shipping = 0 := by
  simp [chargeShipping]
\end{minted}
\caption{simp で分岐を整理する}
\label{lst:ch02_simp}
\end{listing}

コード\ref{lst:ch02_simp}では、\texttt{isMember} が \texttt{true} の場合、\texttt{if true then 0 else shipping} は \texttt{0} に整理される。
\texttt{simp [chargeShipping]} は、\texttt{chargeShipping} の定義を展開し、明らかな分岐を整理している。

\subsection{rw は等式による書き換えである}

\texttt{rw} は rewrite の略である。
すでに分かっている等式を使って、式の一部を書き換える。
リファクタリング証明では非常によく使う。

\begin{listing}[htbp]
\begin{minted}{lean4}
def normalizePrice (price : Nat) : Nat :=
  price

theorem normalizePrice_eq (price : Nat) :
    normalizePrice price = price := by
  rfl

theorem addTax_after_normalize (price tax : Nat) :
    addTax (normalizePrice price) tax = addTax price tax := by
  rw [normalizePrice_eq]
\end{minted}
\caption{rw で既知の等式を使って書き換える}
\label{lst:ch02_rw}
\end{listing}

コード\ref{lst:ch02_rw}では、\texttt{normalizePrice\_eq} という等式を使って、\texttt{normalizePrice price} を \texttt{price} に書き換えている。
その結果、左右が同じ形になり、証明が完了する。

\begin{leanbox}
\texttt{rfl} は「定義を見れば同じ」、\texttt{simp} は「明らかな整理をまとめて行う」、\texttt{rw} は「既知の等式で置き換える」と読む。
この三つは、第3章以降の等式推論とリファクタリング証明の基本語彙になる。
\end{leanbox}

\section{サンプル：税込金額計算の単純な仕様}

\subsection{実務例を小さくする}

ここでは、税込金額計算を題材にする。
現実の税計算には、税率、端数処理、軽減税率、通貨単位、割引との順序などが関わる。
しかし本章では、最初の証明に集中するため、税率ではなく税額を直接受け取る単純なモデルにする。

\begin{softwarebox}
モデルを単純化することは、ごまかしではない。
証明したい核を小さく切り出すための設計判断である。
現実の複雑さをすべて入れる前に、まず「税額を足す」という最小仕様を Lean で確認する。
\end{softwarebox}

\subsection{関数、仕様、証明を並べる}

次のコードは、本章で学んだ要素をまとめたものである。
関数を定義し、命題として仕様を書き、名前付き定理として証明する。

\begin{listing}[htbp]
\begin{minted}{lean4}
def addTax (price tax : Nat) : Nat :=
  price + tax

def TaxSpec (price tax total : Nat) : Prop :=
  total = addTax price tax

example : TaxSpec 1000 100 1100 := by
  rfl

theorem addTax_zero (price : Nat) :
    addTax price 0 = price := by
  simp [addTax]
\end{minted}
\caption{税込金額計算の小さな仕様証明}
\label{lst:ch02_tax_spec}
\end{listing}

コード\ref{lst:ch02_tax_spec}で重要なのは、\texttt{example} と \texttt{theorem} の違いである。
\texttt{example} は、具体的な入力 \texttt{1000} と \texttt{100} についての確認である。
一方、\texttt{addTax\_zero} は任意の \texttt{price} について成り立つ主張である。
テストが具体例を確認するのに対し、定理はすべての入力に対する主張を扱う。

\subsection{ID 保存仕様を書く}

税込金額のような数値例だけでなく、データ変換の仕様も命題として書ける。
次の例では、入力用ユーザーデータを DTO に変換しても、ID が保存されることを証明する。

\begin{listing}[htbp]
\begin{minted}{lean4}
structure UserInput where
  id : Nat
  nameCode : Nat

structure UserDto where
  id : Nat
  displayCode : Nat

def toDto (u : UserInput) : UserDto :=
  { id := u.id, displayCode := u.nameCode }

def PreservesId (f : UserInput -> UserDto) : Prop :=
  forall u, (f u).id = u.id

theorem toDto_preservesId : PreservesId toDto := by
  intro u
  rfl
\end{minted}
\caption{DTO 変換で ID が保存されることを証明する}
\label{lst:ch02_id_preservation}
\end{listing}

コード\ref{lst:ch02_id_preservation}では、\texttt{PreservesId} が仕様である。
\texttt{forall u, ...} は「すべての \texttt{u} について」という意味である。
\texttt{intro u} は、任意のユーザー \texttt{u} を一つ取って証明を始める命令である。
その後、\texttt{toDto} の定義を見れば ID がそのままコピーされているので、\texttt{rfl} で証明できる。

\begin{intuitionbox}
「すべての入力で成り立つ」という仕様は、Lean では \texttt{forall} を使って表す。
証明では、任意の入力を一つ取り出し、その入力について成り立つことを示す。
この形は、後のマイグレーション証明で何度も出てくる。
\end{intuitionbox}

\FigurePlaceholder{fig-ch02-spec-proof-flow.png}{0.92}{自然言語仕様を、関数・命題・証明へ分解する流れ}{fig:ch02-spec-proof-flow}

\section{実務への読み替え}

本章で扱った Lean コードは小さい。
しかし、考え方は実務にそのまま接続できる。

\subsection{テストではなく証明に向く性質}

テストは重要である。
しかし、テストは基本的に選んだ入力例に対する確認である。
一方、Lean の定理は、命題に含まれるすべての入力を対象にする。
この違いにより、次のような性質は証明に向いている。

\begin{itemize}
\item 変換後も ID が保存される。
\item 0 を足しても結果が変わらない。
\item 会員フラグが true のとき送料は 0 になる。
\item 正規化関数が実質的に恒等関数なら、計算結果は変わらない。
\end{itemize}

ただし、すべてを証明に回す必要はない。
外部 API の応答、DB の実際の制約、パフォーマンス、UI の挙動などは、通常のテストや監視の方が適している場合も多い。

\subsection{仕様を小さく名付ける}

Lean で有用なのは、仕様に名前を付けられることである。
たとえば \texttt{toDto\_preservesId} という名前は、何を保証しているかを明確に示している。
この名前を設計文書やレビューコメントに対応させると、証明ファイルが仕様書の索引としても機能する。

\begin{softwarebox}
実務で theorem 名を付けるときは、関数名と保存したい性質を含めると読みやすい。
たとえば \texttt{migrate\_preservesId}、\texttt{rollback\_after\_migrate}、\texttt{map\_migrate\_preservesLength} のような名前である。
\end{softwarebox}

\section{よくある誤解}

\subsection{Bool を返せば仕様を書いたことになる、とは限らない}

\texttt{Bool} を返す関数は、判定を実装している。
しかし、その判定が仕様どおりであることは別の主張である。
たとえば、\texttt{isValidUser} という関数があっても、それが本当にユーザー仕様を満たしているかは別途確認する必要がある。
Lean では、その確認を \texttt{Prop} と theorem で表せる。

\subsection{具体例の証明と一般的な証明は違う}

\texttt{example : TaxSpec 1000 100 1100} は便利である。
しかし、それは特定の値についての確認である。
任意の価格についての仕様を表すには、\texttt{theorem addTax\_zero (price : Nat) : ...} のように、入力を変数として受け取る必要がある。

\subsection{simp は魔法ではない}

\texttt{simp} は強力だが、何でも証明してくれる道具ではない。
定義の展開、既知の単純化規則、明らかな分岐整理などが得意である。
一方、業務仕様そのものを Lean が自動で発見するわけではない。
何を証明したいかは、人間が命題として書く必要がある。

\begin{warningbox}
証明が短いからといって、仕様が軽いわけではない。
よい定義を置いた結果、証明が \texttt{rfl} や \texttt{simp} で終わることはよくある。
それは、仕様と実装の形がうまく揃っているという良いサインである。
\end{warningbox}

\section{本章のまとめ}

\begin{summarybox}
\begin{itemize}
\item \texttt{Bool} は実行時に計算される真偽値であり、\texttt{Prop} は証明したい命題である。
\item 仕様は、関数の外側に命題として切り出すと Lean で検査できる。
\item \texttt{example} はその場限りの確認、\texttt{theorem} は名前付きで再利用できる仕様である。
\item \texttt{by} は証明を書く入口であり、\texttt{rfl}、\texttt{simp}、\texttt{rw} は初期章で使う基本的な証明道具である。
\item 税込金額計算や DTO 変換の ID 保存のような小さな性質は、Lean の命題と証明として表せる。
\end{itemize}
\end{summarybox}

\section{次章への接続}

本章では、関数に対して仕様を命題として書き、その命題に証明を与える方法を学んだ。
次章では、この考え方をさらに進めて、リファクタリングを等式推論として扱う。
変更前後の関数が同じ結果を返すことを、\texttt{rw} や \texttt{calc} を使って段階的に示す。
これは、実務のコード変更を「意味を保存する書き換え」として扱うための最初の本格的なステップになる。
